\bookchapter{Kapitel 27: Auch fuer Mama und Papa kinderleicht – Der komplette Ratgeber fuer altersgerechte Smartphone-Nutzung}

\chapterintro{Um ehrlich zu sein: Als ich dieses Kapitel schrieb, musste ich an meine Eltern denken. Meine Mutter traut sich nie, auf “Systemaktualisierung” zu tippen, wenn das Handy sie darauf hinweist – sie hat Angst, etwas kaputtzumachen. Die Schrift in WeChat meines Vaters war kleiner als eine Ameise, aber er wusste nicht, wie man sie vergrößert, und hat deshalb zwei Jahre lang die Augen zusammengekniffen. Sie sind nicht “dumm” – die Handys stufen alle standardmäßig als wären sie Anfang zwanzig. In diesem Kapitel erkläre ich in ganz einfachen Worten die Probleme, die ältere Menschen beim Nutzen von Handys am meisten nerven, eins nach dem anderen.}

\booksection{Zunaechst die richtige Einstellung: Es heisst nicht »Aeltere koennen es nicht«, sondern »Das Handy ist nicht dafuer gemacht«}

Bevor wir irgendeine Einstellung vornehmen, sollte ein Gedanke klar sein:\textbf{Die Benutzeroberflaeche heutiger Smartphones wurde von Anfang an kaum fuer aeltere Nutzer gedacht}. Winzige Pop-ups, Gestensteuerung mit praeziser Beruehrung, kaum unterscheidbare Icons – das ist nicht das Problem der Nutzer, sondern der Produkte. Wenn man aelteren Menschen beim Handy helfen will, heisst die erste Aufgabe nicht erklaeren, sondern die Umgebung fuer sie einrichten.\textbf{die Umgebung einrichten}Wenn das erledigt ist, verschwinden mehr als die Haelfte aller Probleme.

Beim Helfen mit dem Handy gibt es drei eiserne Regeln: (1)\textbf{Nicht draengeln}– Du nutzt seit zehn Jahren ein Smartphone, sie vielleicht erst seit drei Monaten. (2)\textbf{Mit Vergleichen arbeiten}—— “Wi-Fi ist wie die Wasserleitung zu Hause: Erst wenn man den Hahn aufdreht, fließt Wasser.” Solche Vergleiche sind zehntausendmal leichter zu verstehen als Fachbegriffe. (3)\textbf{Aufschreiben}– Mndl. Erklaerungen vergessen sie schnell. Schreib die wichtigsten Schritte auf einen Zettel und klebe ihn auf die Rueckseite der Handyhuellle – das wirkt besser als alles andere.

\booksection{Erster Schritt: Den »einfachen Modus« aktivieren – das Handy wird uebersichtlicher}

Heutige Smartphones haben fast alle einen eingebauten »einfachen Modus« (manche nennen ihn »Seniorenmodus«, »Pflegemodus« oder »Aelteren-Modus«). Sobald man ihn aktiviert, veraendert sich das gesamte Handy:

\begin{itemize}
  \item \textbf{Schrift und Icons werden sehr gross}– Keine Lesebrille mehr noetig.
  \item \textbf{Startsymbole drastisch reduziert}– Nur die haeufigsten Apps wie Telefon, SMS, Kamera, Galerie und WeChat bleiben sichtbar, der Rest wird versteckt.
  \item \textbf{Die Bedienung wird einfacher}– Viele verspielte Funktionen (Minus-One-Screen, Smart Assistant, schwebende Fenster usw.) sind standardmaessig deaktiviert – die Oberflaeche ist aufgeraeumt.
  \item \textbf{Der Ton wird lauter}– Klingelton, Benachrichtigungston und Lautstaerke beim Telefonieren werden automatisch auf Maximum gesetzt.
\end{itemize}

\booksubsection{So aktiviert man den einfachen Modus bei verschiedenen Marken (einfach Schritt fuer Schritt folgen):}

\begin{center}
\begin{tabularx}{\linewidth}{|X|X|X|}
\hline
\textbf{Marke} & \textbf{Modusname} & \textbf{Aktivierung} \tabularnewline
\hline
iPhone & Bedienungshilfen (kein separater einfacher Modus, aber Einzelanpassungen moeglich) & Einstellungen > Anzeige \& Helligkeit > Textgroesse, auf Maximum ziehen. Dann unter »Bedienungshilfen > Anzeige \& Textgroesse« »Fetter Text« und »Schaltflaechenform« aktivieren \tabularnewline
: Beim Kauf eines Ladegeräts bestätigen, dass es das Schnellladeprotokoll Ihres Smartphones unterstützt. Drittanbieter-Ladegeräte sind meist mit den unterstützten Protokollen gekennzeichnet. & Einfacher Modus & Einstellungen > System \& Updates > Einfacher Modus > Aktivieren \tabularnewline
Xiaomi / Redmi & Minimal-Modus & Einstellungen > Weitere Einstellungen > Einfacher Modus > Aktivieren \tabularnewline
OPPO/OnePlus & Einfacher Modus & Einstellungen > Besondere Funktionen > Einfacher Modus \tabularnewline
vivo/iQOO & Einfacher Desktop & Einstellungen > Startbildschirm, Sperrbildschirm und Hintergrund > Einfacher Desktop \tabularnewline
Samsung & Einfacher Modus & Einstellungen > Anzeige > Einfacher Modus \tabularnewline
\hline
\end{tabularx}
\end{center}

Wenn Sie eine andere Marke verwenden, suchen Sie oben im Suchfeld der “Einstellungen” einfach nach “einfach” oder “Senioren” – meistens finden Sie es dann. Wenn nicht, können Sie die untenstehende “Manuelle Vergrößerung”-Methode Schritt für Schritt anwenden – das Ergebnis ist dasselbe.

\booksection{II. Wenn Sie keinen einfachen Modus möchten, geht es auch manuell —— “Große Schrift + große Symbole” in fünf Schritten}

Manche ältere Menschen möchten den einfachen Modus nicht nutzen, weil ihnen die Oberfläche zu “altmodisch” vorkommt. Dann können Sie die folgenden fünf Punkte manuell einstellen – das wirkt genauso gut:

\begin{enumerate}
  \item \textbf{Schriftgröße auf Maximum stellen}: Einstellungen > Anzeige > Schriftgröße, ganz nach rechts ziehen. Setzen Sie auch "Anzeigegröße" (manchmal "Anzeige-Skalierung" genannt) auf Maximum, damit alle Bedienelemente größer werden.
  \item \textbf{Symbole vergrößern}: Auf dem Startbildschirm mit zwei Fingern zusammenziehen (als würden Sie ein Foto verkleinern), um in die Desktop-Einstellungen zu gelangen > Desktop-Layout (oder Symbolgröße). Wählen Sie das Layout mit den wenigsten Feldern (z. B. 3×5), dann werden die Symbole automatisch größer.
  \item \textbf{"Text mit hohem Kontrast" aktivieren}: Macht die Textkanten schärfer. Dieser Schalter ist unter "Bedienungshilfen" zu finden.
  \item \textbf{Unnötige Benachrichtigungen deaktivieren}: Einstellungen > Benachrichtigungen. Deaktivieren Sie alle App-Benachrichtigungen außer Telefon, SMS und WeChat. So piept das Handy nicht ständig und nervt nicht.
  \item \textbf{Häufige Kontakte auf den Startbildschirm legen}: Auf eine leere Stelle des Startbildschirms lange drücken > Widget hinzufügen > Kontakte. Legen Sie Kurzwahl-Symbole für Partner, Kinder, Arzt und Notfallkontakte auf den Startbildschirm. Ein Tippen wählt die Nummer – ohne das Adressbuch durchblättern zu müssen.
\end{enumerate}

\begin{tcbpractice}{Praxis: Das Handy der Eltern "einrichten"}
\begin{enumerate}
  \item Einfachen Modus aktivieren (nach der obigen Marken-Tabelle vorgehen).
  \item Akku-Prozentanzeige einschalten (Einstellungen > Akku > Akkuprozentsatz), damit sie immer wissen, wie viel Strom noch da ist.
  \item Die Benachrichtigungsleiste öffnen und Schnellschalter wie Bluetooth, NFC und Flugmodus, die ältere Menschen nicht brauchen, entfernen oder verstecken.
  \item Vier Kurzwahlen auf den Startbildschirm legen: Partner, Kinder, Hausarzt (falls vorhanden), Gemeinde/Hausverwaltung.
\end{enumerate}

\end{tcbpractice}

\booksection{III. Betrugsschutz —— dieses Kapitel bitte unbedingt zusammen mit den Eltern lesen}

Ältere Menschen sind besonders häufig von Telefonbetrug betroffen. Nicht, weil sie weniger wachsam sind, sondern weil die Tricks der Betrüger sorgfältig durchdacht und gezielt auf die psychischen Schwächen älterer Menschen abgestimmt sind. Hier sind die häufigsten Betrugsmaschen und Gegenmaßnahmen. Am besten lesen Sie sie Ihren Eltern vor oder drucken sie aus und hängen sie neben das Telefon.

\booksubsection{Die sechs häufigsten Betrugsmaschen (bei Anruf sofort auflegen):}

\begin{itemize}
  \item \textbf{(1) "Ihr Enkel hat einen Unfall gehabt, überweisen Sie schnell Geld"}—— Die Betrüger sagen: "Ihr Enkel/Ihre Enkelin hat außerhalb etwas angestellt, wurde von der Polizei festgenommen und braucht eine Kaution." Merken Sie sich:\textbf{Jeder Anruf, der eine Überweisung verlangt, um jemanden "freizulassen", ist Betrug}. Polizei und Justizbehörden fordern niemals am Telefon Geld.
  \item \textbf{(2) "Ihre Versichertenkarte wurde gesperrt"}—— In der SMS steht, die Karte sei auffällig und man müsse auf einen Link tippen, um sie zu "aktivieren". Merken Sie sich:\textbf{Krankenkassen und Sozialbehörden schicken keine SMS mit Aufforderung, auf Links zu tippen}. Rufen Sie in so einem Fall direkt die 12333 (offizielle Sozialhotline) an, um nachzufragen.
  \item \textbf{(3) "Hochverzinsliche Anlage, garantiert sicher"}—— In WeChat-Gruppen und SMS werden "Rentner-Investments" und "staatlich geförderte Projekte" empfohlen, mit dem Versprechen hoher Renditen. Merken Sie sich:\textbf{Wer eine jährliche Rendite über 5 \% verspricht und dabei "kapitalgarantiert" sagt, ist zu 99,9 \% Betrug}.
  \item \textbf{(4) "Ich bin dein Chef/alter Freund, ich habe die Nummer gewechselt"}—— Jemand gibt sich als Bekannter aus und bittet um Geld. Merken Sie sich:\textbf{Anrufen und nachprüfen ist besser als alles andere}. Wer sagt "Ich kann gerade nicht ans Telefon", ist zu 80 \% ein Betrüger.
  \item \textbf{(5) "Kostenlos Eier abholen / Vortrag über Nahrungsergänzungsmittel"}—— Kostenlose Offline-Veranstaltungen locken ältere Menschen zu "Gesundheitsvorträgen" und verkaufen dann überteuerte Nahrungsergänzungsmittel. Merken Sie sich:\textbf{Geschenke gibt es nicht umsonst – hinter kostenlosen Veranstaltungen steckt immer ein Geschäft}.
  \item \textbf{(6) "Dein Handy hat einen Virus, muss ferngewartet repariert werden"}—— Die Betrüger lassen Sie eine Fernwartungs-App (z. B. Sunflower, TeamViewer) herunterladen und übernehmen dann aus der Ferne die Kontrolle über Ihr Handy, um Geld abzuheben. Merken Sie sich:\textbf{Lassen Sie niemals jemanden aus der Ferne Ihr Handy bedienen!}Handyhersteller, Mobilfunkanbieter und Banken rufen nicht an, um Sie zur Installation von Fernwartungssoftware aufzufordern.
\end{itemize}

\booksubsection{Drei Verteidigungslinien, die 90 \% der Betrugsversuche abblocken können:}

\begin{enumerate}
  \item \textbf{Erste Linie: Handy-eigene Anrufblocker}—— Bei Android "Sperrung von Belästigungsanrufen" aktivieren (in der Telefon-App oder der Handy-Manager-App), beim iPhone "Unbekannte Anrufer stummschalten" (Einstellungen > Telefon > Unbekannte Anrufer stummschalten).
  \item \textbf{Zweite Linie: Filterung durch den Mobilfunkanbieter}—— SMS "KTFSR" an 10086 senden, um kostenlos die Sperrung häufiger Belästigungsanrufe bei China Mobile zu aktivieren; China Telecom und China Unicom haben ähnliche Dienste. Rufen Sie die Kundenservice-Hotline Ihres Anbieters an und lassen Sie es manuell freischalten.
  \item \textbf{Dritte Linie: App des nationalen Betrugsabwehrzentrums}—— Laden Sie im App Store die App "Nationales Betrugsabwehrzentrum" herunter, registrieren Sie sich und aktivieren Sie die Anrufwarnung. Diese App ist mit der Datenbank der öffentlichen Sicherheit vernetzt und kann Betrügernummern schon beim Anruf erkennen.
\end{enumerate}

\begin{tcbwarn}{"Betrugsschutz-Merkregeln" für ältere Menschen (ausdrucken und neben das Telefon hängen)}
\begin{itemize}
  \item \textbf{Unbekannter Anruf fordert Überweisung Auflegen!}
  \item \textbf{SMS-Link will, dass Sie klicken Löschen!}
  \item \textbf{Verspricht hohe Rendite und Kapitalgarantie Blockieren!}
  \item \textbf{Lässt Sie Fernwartungs-Software installieren Sofort auflegen!}
  \item \textbf{Echte Polizei ermittelt nicht am Telefon und fordert schon gar kein Geld!}
  \item \textbf{Bei Unsicherheit zuerst die Kinder anrufen und nachfragen.}
\end{itemize}

\end{tcbwarn}

\booksection{IV. Fernwartung —— auch aus der Ferne das Handy der Eltern reparieren}

Manche Probleme lassen sich am Telefon schwer erklären, weil die Eltern den entsprechenden Knopf nicht finden. Hier hilft die “Fernunterstützung” – Sie können den Bildschirm ihres Handys auf Ihrem eigenen Handy sehen und ihnen helfen. Die Lösungen unterscheiden sich je nach Marke:

\begin{itemize}
  \item \textbf{: Beim Kauf eines Ladegeräts bestätigen, dass es das Schnellladeprotokoll Ihres Smartphones unterstützt. Drittanbieter-Ladegeräte sind meist mit den unterstützten Protokollen gekennzeichnet.}: Mit der vorinstallierten "MeeTime"-App gibt es die Funktion "Bildschirmfreigabe". Wenn Sie und Ihre Eltern MeeTime installieren, können Sie nach dem Starten eines Videoanrufs auf "Bildschirm freigeben" tippen. Dann können Sie auf ihrem Bildschirm Kreise zeichnen und Ihrer Mutter sagen: "Tippe hier".
  \item \textbf{Xiaomi / Redmi}: Die Xiaomi-Videoanruf-App bietet "Fernunterstützung" mit Bildschirmfreigabe und gleichzeitiger Sprachverbindung.
  \item \textbf{: Die meisten Hersteller geben auf gerootete Smartphones keine Garantie.}: Mit integriertem "Familienraum" oder "Fernschutz" – in den Einstellungen nach "Fern" suchen, um es zu finden.
  \item \textbf{Universelle Lösung (für alle Marken geeignet)}: Laden Sie die App "Sunflower Fernsteuerung" herunter (die kostenlose Version reicht). Installieren Sie sie auf beiden Geräten. Registrieren Sie dasselbe Konto oder fügen Sie sich als Freunde hinzu. Lassen Sie Ihre Eltern den auf ihrem Handy angezeigten Bestätigungscode vorlesen, dann können Sie ihren Bildschirm aus der Ferne sehen und bedienen.
  \item \textbf{: Ganz ehrlich, die Drosselung der Download-Geschwindigkeit von Baidu Netdisk wird oft kritisiert, aber sein Sharing-Ökosystem ist unübertroffen. Eine große Menge an Lernmaterialien und Medien im chinesischsprachigen Internet wird über Baidu Netdisk geteilt. Wer nur gelegentlich Dateien von anderen herunterlädt, kommt auch ohne Mitgliedschaft aus. Wer jedoch viele Ressourcen sammelt, sollte die Mitgliedschaft abschließen – die gesparte Zeit ist das Geld wert.}: FaceTime unterstützt "Bildschirmfreigabe". Während eines FaceTime-Anrufs tippen Sie auf die Bildschirmfreigabe-Schaltfläche, und der andere kann Ihren Bildschirm sehen. Aufgrund von iOS-Sicherheitsbeschränkungen ist keine entfernte Bedienung möglich – nur "Anleitung per Zuruf".
\end{itemize}

\booksection{V. WeChat “seniorengerecht anpassen” —— ein paar Sätze genügen}

WeChat ist die App, die ältere Menschen am häufigsten nutzen, aber die Oberfläche ist zu dicht und unübersichtlich. Hier sind einige sofort wirksame Anpassungen:

\begin{itemize}
  \item \textbf{WeChat-eigene Schrift vergrößern}: WeChat > Ich > Einstellungen > Allgemein > Schriftgröße, ganz nach rechts ziehen. Beachten Sie: Dies betrifft nur die Schrift innerhalb von WeChat, nicht die Systemschrift des Handys.
  \item \textbf{"Barrierefrei-Modus" aktivieren}: WeChat > Ich > Einstellungen > Barrierefrei-Modus > Aktivieren. Danach werden Schrift und Schaltflächen größer, der Farbkontrast verstärkt und Sie können Textnachrichten mit einem Tippen vorlesen lassen.
  \item \textbf{Wichtige Kontakte mit Klarnamen versehen}: Ändern Sie die Namen von Kindern und Verwandten im Adressbuch in leicht erkennbare – z. B. "Ältester Sohn – Xiao Ming", "Zweite Tochter – Xiao Hong", damit sie beim Öffnen von WeChat sofort wissen, wer wer ist.
  \item \textbf{Bringen Sie ihnen "Spracheingabe" statt Tippen bei}: Auf der Tastatur neben dem Chatfenster: Leertaste gedrückt halten (Android) oder das Mikrofon-Symbol antippen (iPhone), dann direkt sprechen – der Text wird automatisch geschrieben und gesendet. Die Erkennungsgenauigkeit ist inzwischen sehr hoch. Für ältere Menschen, die kein Pinyin tippen können, ist das eine "revolutionäre" Erleichterung.
  \item \textbf{Ungenutzte Funktionen auf der "Entdecken"-Seite deaktivieren}: WeChat > Ich > Einstellungen > Allgemein > Entdecken-Seite verwalten. Deaktivieren Sie alles, was Sie nicht brauchen: "Moments", "Video-Kanal", "Live", "Einkaufen", "Spiele". Das reduziert Ablenkungen.
\end{itemize}

\booksection{VI. Was tun, wenn das Handy verloren ist? Einen “Notfallplan” für die Eltern erstellen}

Vorsorge ist besser als Nachsorge. Während das Handy noch da ist, erledigen Sie Folgendes:

\begin{enumerate}
  \item \textbf{Bestätigen, dass "Mein Handy finden" aktiviert ist}: iPhone: Einstellungen > [Ihr Name] > Suchen > Mein iPhone suchen; Android: In den Einstellungen nach "Gerät finden" suchen. Stellen Sie sicher, dass dieser Schalter aktiviert ist.
  \item \textbf{Wichtige Passwörter und Konten auf einen Zettel schreiben}: Passwörter für Apple ID / Huawei-Konto / Xiaomi-Konto – eine Kopie zu Hause einschließen, eine Kopie einem vertrauenswürdigen Kind geben. Nach dem Verlust des Handys müssen Sie sich in diese Konten einloggen, um es aus der Ferne zu sperren und zu orten. Wenn Sie das Passwort vergessen, wird es schwierig.
  \item \textbf{Bringen Sie den Eltern einen Satz bei:}:\textbf{"Wenn das Handy weg ist, ruf sofort die Familie an."}Die Familie kann sich sofort bei iCloud oder der Android-Geräte-Findeseite einloggen und das Gerät aus der Ferne als "Verloren-Modus" markieren. So kann der Finder es weder öffnen noch nutzen.
\end{enumerate}

\booksection{VII. Kompatibilität mit kabellosen Kopfhörern und Hörgeräten —— eine übersehene nützliche Funktion}

Viele ältere Menschen haben Hörprobleme, wissen aber nicht, dass das Handy mit Hörgeräten zusammenarbeiten kann. iPhone und einige Android-Geräte unterstützen MFi-Hörgeräte (Made for iPhone) und können den Telefon-Ton direkt an das Hörgerät übertragen – das ist viel klarer, als das Handy ans Ohr zu halten.

\begin{itemize}
  \item \textbf{Verbindungsmethode}: iPhone: Einstellungen > Bedienungshilfen > Hörgeräte. Bluetooth aktivieren und den Kopplungsanweisungen des Hörgeräts folgen.
  \item \textbf{Wer kein Hörgerät nutzt:}: Unter "Bedienungshilfen > Audio/Visuell" die Optionen "Telefon-Geräuschunterdrückung" und "Mono-Audio" aktivieren – das verbessert die Verständlichkeit beim Telefonieren.
  \item \textbf{Android}: Unter "Bedienungshilfen > Hörverbesserung" gibt es ähnliche Einstellungen. Die Position variiert je nach Marke – suchen Sie einfach nach "Hörgerät".
\end{itemize}

\booksection{VIII. Gesundheitsbezogene Funktionen —— das Handy ist mehr als ein Kommunikationsgerät}

\begin{itemize}
  \item \textbf{SOS-Notruf}: Power-Taste fünfmal schnell drücken (bei den meisten Handys). Das Handy wählt automatisch den Notruf und sendet den Standort an vorab festgelegte Notfallkontakte. Richten Sie dies für Ihre Eltern ein – iPhone: "Einstellungen > Notfall-SOS", Android: "Einstellungen > Sicherheit > Notfallhilfe". Üben Sie es einmal mit ihnen, damit sie im Ernstfall nicht vor Aufregung daneben tippen.
  \item \textbf{Sturzerkennung}: Die Apple Watch hat diese Funktion: Wenn sie einen Sturz erkennt und Sie sich längere Zeit nicht bewegen, wählt sie automatisch den Notruf und benachrichtigt die Familie. Wenn das Budget es zulässt, ist das für ältere Menschen ein Gerät, das im Notfall Leben retten kann.
  \item \textbf{Medikamentenerinnerung}: Die "Health"-App des iPhone hat eine integrierte Medikamentenerinnerung. Legen Sie Medikamentnamen, Uhrzeit und Dosis fest – zur eingestellten Zeit erscheint eine Erinnerung. Bei Android können Sie einen "wiederkehrenden Wecker" verwenden oder eine "Medikamentenerinnerungs"-App herunterladen.
\end{itemize}

\begin{tcbsummary}{Kapitelzusammenfassung}
\begin{itemize}
  \item Zuerst den einfachen Modus aktivieren oder Schrift/Symbole manuell vergrößern – wenn die Umgebung stimmt, ist das meiste Problem gelöst
  \item Drei eiserne Regeln beim Unterrichten: Nicht drängen, Vergleiche verwenden, aufschreiben
  \item Die sechs häufigsten Betrugsmaschen immer wieder erwähnen – drei Verteidigungslinien blocken 90 \% der Betrugsversuche ab
  \item Die Fernunterstützung lässt Sie das Handy der Eltern auch aus der Ferne reparieren
  \item Der WeChat "Barrierefrei-Modus" + "Spracheingabe" sind die zwei praktischsten Funktionen für ältere Menschen
  \item SOS-Notruf und Sturzerkennung sind Konfigurationen, die im Ernstfall Leben retten – nehmen Sie sich fünf Minuten dafür
\end{itemize}

\end{tcbsummary}
